\section{Related work}
\label{sec:related}

\paragraph{Counterfactual explanations.} \citet{wachter2017counterfactual}
framed a counterfactual as the closest input that changes the decision, found by
minimising a distance-to-target objective. The idea was extended to produce
\emph{diverse} sets \citep{mothilal2020dice}, to stay on the data manifold with
prototypes \citep{vanlooveren2021interpretable} or density-weighted paths
\citep{poyiadzi2020face}, and to trade objectives explicitly with multi-objective
search \citep{dandl2020multi}. Surveys catalogue dozens of such methods
\citep{verma2020counterfactual,karimi2022survey} and benchmarks compare them
\citep{pawelczyk2021carla,klaise2021alibi}. The near-universal headline metric is
validity, sometimes with proximity and diversity; the question of whether the
named point could be \emph{reached} is largely left aside, and the assumptions
that licence reading a point as advice --- above all that its features can be set
independently of one another --- typically go unstated \citep{barocas2020hidden}.

\paragraph{Actionable recourse.} \citet{ustun2019actionable} made mutability
first-class: some features cannot be changed, others only in one direction, and
recourse must live in the actionable subspace. This is an essential correction,
and we adopt it wholesale as one of our three audits. It is, however, a constraint
on the \emph{named} change; on its own it still treats the world as a set of
independently-settable dials.

\paragraph{Causal recourse.} \citet{karimi2021recourse} argued that recourse is
about \emph{interventions}, not points: because features are causally linked, one
should optimise over actions $\dopart{X_S = v_S}$ and read off consequences
through a structural causal model \citep{pearl2009causality}, and they show a
minimal-cost intervention can differ sharply from the nearest counterfactual.
\citet{karimi2020imperfect} extend this to imperfect causal knowledge. Our method
sits squarely in this line --- it is the mechanism by which \cac{} is causally
valid by construction. Our contribution relative to it is (i) an \emph{evaluation}
that scores existing point-based methods on causal validity and actionability
together, quantifying how far short they fall, and (ii) a single search that adds
an explicit human-effort cost, an actionability guard, and a plausibility term
around the causal core, rather than causal validity alone.

\paragraph{Robustness.} A separate concern is whether recourse survives model
change \citep{upadhyay2021robust}; we touch it only as one axis of our robustness
study (\S\ref{sec:experiments}), where we re-audit \cac{} under a perturbed SCM, a
swapped black box, a tighter budget, and a different plausibility neighbourhood.

\paragraph{Position.} The pieces --- validity, actionability, causal consequence
--- exist. What has been missing, and what we supply, is (a) an evaluation that
insists on \emph{all three at once} and reports their conjunction, and (b) a
method that treats causal validity as a property of the search space rather than a
soft penalty, so it cannot be traded away for a lower distance.
